\chapter{Theoretical Background}

This chapter presents the model and the tools used in the rest of the work: the Hawkes
process, from the univariate case to the symmetric bivariate form used here, its stability
and the volatility it implies; simulation, maximum likelihood, model comparison and the
goodness-of-fit test built on the compensator; the effect of observing the events on a time
grid; the reading of a network trained on simulated data as a Bayes estimator; the statistical
tools used to compare estimators; and the two architectures compared, with their training.

\section{Point processes}

If \(t_1<t_2<\cdots\) are event times, the counting process is
\[
  N(t)=\sum_k\mathbf{1}_{\{t_k\le t\}}.
\]
The conditional intensity \(\lambda(t\mid\mathcal H_t)\) describes the expected
instantaneous rate given the history \(\mathcal H_t\) \citep{daley2003}. Informally,
\[
  \Pr[N(t+\mathrm dt)-N(t)=1\mid\mathcal H_t]
  \approx \lambda(t)\,\mathrm dt.
\]
It is not a probability in arbitrary units: intensity has the dimension of events per
unit of time. All processes in this work are simple: two events never occur at the same
instant, so the probability of more than one event in \(\mathrm dt\) is of smaller order.
The intensity then determines the law of the process. Its integral,
\[
  \Lambda(t)=\int_0^t\lambda(u)\,\mathrm du,
\]
is the compensator: the number of events expected up to \(t\) given the past, so that
\(N(t)-\Lambda(t)\) is a martingale. The compensator appears twice below, in the likelihood
(Section~\ref{sec:mle-theory}) and in the goodness-of-fit test
(Section~\ref{sec:time-rescaling}).

\section{Poisson as the null model}

In the homogeneous Poisson process, \(\lambda(t)=\lambda_0\). Inter-event times are
independent exponentials. The model produces no bursts beyond those due to chance. It
is therefore a useful control: if Hawkes does not outperform Poisson in diagnostics,
self-excitation is not justified. In the inhomogeneous Poisson process the intensity is a
deterministic function \(\lambda_0(t)\), so it can follow, for example, an intraday pattern of
activity, but it still ignores the past. The Hawkes process keeps the dependence on the past
and lets each event change the rate of the events that follow.

\section{Univariate Hawkes process}

With an exponential kernel,
\begin{equation}
  \lambda(t)=\mu+\sum_{t_k<t}\alpha e^{-\beta(t-t_k)}.
  \label{eq:hawkes1d}
\end{equation}
\(\mu\) is the exogenous activity; \(\alpha\) is the immediate jump; \(\beta\)
controls the speed of forgetting. The integral of the kernel is \(\alpha/\beta\), the
branching ratio \citep{laub2021}. Each event raises the intensity by \(\alpha\), and the
increase decays with half-life \(\ln2/\beta\) (Figure~\ref{fig:intensity}).

In general, \(\lambda(t)=\mu+\int_{(0,t)}\phi(t-s)\,\mathrm dN(s)\) for a non-negative kernel
\(\phi\) \citep{hawkes1971}. The exponential kernel \(\phi(s)=\alpha e^{-\beta s}\) has a single
time scale. A sum of exponentials, \(\sum_m\alpha_m e^{-\beta_m s}\), mixes several scales and
keeps the recursion below for each term. A power law, \(\phi(s)\propto(c+s)^{-(1+\gamma)}\),
has long memory, and fits of this kind to financial data tend to place the process close to
criticality \citep{hardiman2013,bacry2015}. The choice matters for the diagnostics: a single
exponential fitted to data with several time scales leaves dependence in the residuals
(Section~\ref{sec:time-rescaling}).

\begin{figure}[htbp]
\centering
\begin{tikzpicture}
\begin{axis}[width=0.9\linewidth,height=5.2cm,xmin=0,xmax=5,ymin=-0.35,ymax=1.75,
  axis lines=left,xlabel={time $t$ (s)},ylabel={$\lambda(t)$},
  ytick={0,0.5,1,1.5},xtick={0,1,1.5,2,3,4,5},
  every axis plot/.append style={line width=1.1pt},clip=false]
  \addplot[name path=lam,color=black,samples=500,domain=0:5]
    {0.5+(x>1)*0.8*exp(-2*max(x-1,0))+(x>1.5)*0.8*exp(-2*max(x-1.5,0))+(x>4)*0.8*exp(-2*max(x-4,0))};
  \addplot[name path=zero,draw=none,domain=0:5]{0};
  \addplot[gray!20] fill between[of=lam and zero];
  \addplot[gray,dashed,thin,domain=0:5]{0.5};
  \node[font=\scriptsize,anchor=south west] at (axis cs:2.9,0.5) {baseline $\mu=0.5$};
  \foreach \t/\l in {1/0.5,1.5/0.794,4/0.507}{
    \edef\temp{\noexpand\draw[line width=1.1pt] (axis cs:\t,-0.3) -- (axis cs:\t,-0.05);
      \noexpand\fill (axis cs:\t,\l) circle (2pt);}
    \temp}
  \node[font=\scriptsize,anchor=north west] at (axis cs:1.58,0.76) {$\lambda(t_2)=0.794$};
  \draw[{Stealth}-{Stealth},thick] (axis cs:0.93,1.30) -- node[left,font=\scriptsize]{jump $\alpha=0.8$} (axis cs:0.93,0.52);
  \node[font=\scriptsize,align=center] at (axis cs:3.0,0.2) {shaded area: compensator $\Lambda(T)$};
\end{axis}
\end{tikzpicture}
\caption{Intensity of a univariate Hawkes process with \(\mu=0.5\), \(\alpha=0.8\) and
\(\beta=2\), for three events (ticks on the axis). Each event adds a jump that decays with
half-life \(\ln 2/2\approx0.35\)~s. The dots are the intensities at the events, the terms
that the likelihood rewards; the shaded area is the compensator, the term it penalizes.}
\label{fig:intensity}
\end{figure}

\paragraph{The intensity as a state.} Write \(E(t)=\lambda(t)-\mu\) for the excitation.
Between events, \(E\) decays as \(\mathrm dE/\mathrm dt=-\beta E\); at an event it jumps by
\(\alpha\). Hence, with \(\Delta_k=t_k-t_{k-1}\),
\begin{equation}
  E(t_k^-)=e^{-\beta\Delta_k}\bigl(E(t_{k-1}^-)+\alpha\bigr),
  \label{eq:recursion}
\end{equation}
and the current intensity summarizes the whole past: with an exponential kernel the
process \(\lambda(t)\) is Markov. This property makes simulation and the likelihood linear
in the number of events. It does not hold for other kernels, such as the power law, where
the whole history must be kept.

\paragraph{Cluster representation.} The same process can be generated as a population
\citep{hawkes2018,laub2021}. Immigrants arrive as a Poisson process of rate \(\mu\); each
event, immigrant or not, produces children as an inhomogeneous Poisson process of rate
\(\alpha e^{-\beta s}\), where \(s\) is the time since the parent. The number of direct
children of an event is Poisson with mean
\[
  g=\int_0^\infty\alpha e^{-\beta s}\,\mathrm ds=\frac{\alpha}{\beta},
\]
which is why \(g\) is called the branching ratio. Children, grandchildren and further
generations form a geometric series, so the mean size of the cluster started by one
immigrant is
\[
  1+g+g^2+\cdots=\frac{1}{1-g}\qquad(g<1),
\]
and a fraction \(g\) of all events is endogenous, triggered by earlier events rather than by
the baseline. Taking expectations in Equation~\eqref{eq:hawkes1d}, the stationary mean rate
\(\bar\lambda=\E[\lambda(t)]\) satisfies \(\bar\lambda=\mu+g\bar\lambda\), that is,
\(\bar\lambda=\mu/(1-g)\). Only the times are observed; the genealogy
(Figure~\ref{fig:cluster}) is latent.

\begin{figure}[htbp]
\centering
\begin{tikzpicture}[x=1.0cm,y=1.1cm,
  imig/.style={circle,fill=black,inner sep=2.4pt},
  child/.style={circle,draw=black,fill=white,thick,inner sep=2pt},
  link/.style={-{Stealth[length=1.6mm]},gray,thin}]
  \draw[-{Stealth}] (0,0) -- (12.4,0) node[right,font=\small]{$t$};
  \node[imig] (a) at (0.8,0) {};
  \node[child] (a1) at (1.4,1) {}; \node[child] (a2) at (2.3,1) {};
  \node[child] (a11) at (2.0,2) {};
  \draw[link] (a) to[bend left=20] (a1); \draw[link] (a) to[bend left=20] (a2);
  \draw[link] (a1) to[bend left=20] (a11);
  \node[imig] (b) at (4.6,0) {};
  \node[child] (b1) at (5.3,1) {};
  \draw[link] (b) to[bend left=20] (b1);
  \node[imig] (c) at (7.4,0) {};
  \node[imig] (d) at (9.0,0) {};
  \node[child] (d1) at (9.5,1) {}; \node[child] (d2) at (10.2,1) {}; \node[child] (d3) at (11.1,1) {};
  \node[child] (d21) at (10.9,2) {}; \node[child] (d22) at (11.6,2) {};
  \draw[link] (d) to[bend left=20] (d1); \draw[link] (d) to[bend left=20] (d2); \draw[link] (d) to[bend left=20] (d3);
  \draw[link] (d2) to[bend left=20] (d21); \draw[link] (d2) to[bend left=20] (d22);
  \foreach \p in {a,a1,a2,a11,b,b1,c,d,d1,d2,d3,d21,d22}{\draw[densely dotted,gray] (\p) -- (\p|-{(0,-0.3)});}
  \foreach \y in {0,1,2}{\node[font=\scriptsize,anchor=east] at (-0.1,\y) {generation \y};}
  \node[font=\scriptsize,anchor=north west] at (0,-0.4) {$\bullet$ immigrant (rate $\mu$)\qquad
    $\circ$ child (each event has, on average, $g$ children)};
\end{tikzpicture}
\caption{Cluster representation of a Hawkes process. Only the projection of the events on
the time axis is observed; which event triggered which is hidden. With \(g=0.5\), each
immigrant brings on average a cluster of \(1/(1-0.5)=2\) events.}
\label{fig:cluster}
\end{figure}

\section{Symmetric bivariate Hawkes process}

A \(D\)-dimensional Hawkes process has one intensity per type,
\[
 \lambda_i(t)=\mu_i+\sum_{j=1}^{D}\int_{(0,t)}\phi_{ij}(t-s)\,\mathrm dN_j(s),
 \qquad i=1,\dots,D,
\]
where \(\phi_{ij}\) is the effect of a type-\(j\) event on the intensity of type \(i\)
\citep{hawkes1971b,bacry2015}. With exponential kernels and two types there are, in general,
two baselines, four jump sizes and up to four decay rates. Here,
up-moves and down-moves are processes \(N_1,N_2\) with a common decay and symmetric
excitation:
\begin{align}
\lambda_1(t)&=\mu+\sum_{t_k^1<t}\alpha_{\mathrm{self}}e^{-\beta(t-t_k^1)}
 +\sum_{t_k^2<t}\alpha_{\mathrm{cross}}e^{-\beta(t-t_k^2)},\\
\lambda_2(t)&=\mu+\sum_{t_k^2<t}\alpha_{\mathrm{self}}e^{-\beta(t-t_k^2)}
 +\sum_{t_k^1<t}\alpha_{\mathrm{cross}}e^{-\beta(t-t_k^1)}.
\end{align}
The model reduces a general matrix to four parameters. This parsimony makes
simulation and comparison with \citet{lee2023} easier. Figure~\ref{fig:bivariate} draws it
as a diagram of the two event types. An up-move raises the intensity of another up-move by
\(\alpha_{\mathrm{self}}\) and that of a down-move by \(\alpha_{\mathrm{cross}}\), and
symmetrically for a down-move; both increases decay at the same rate \(\beta\). In
microstructure, dominant cross-excitation indicates reversion (an up-move followed by a
down-move), and dominant self-excitation indicates persistence of direction
\citep{bacry2013}.

\begin{figure}[htbp]
\centering
\begin{tikzpicture}[node distance=4.4cm,
  ev/.style={circle,draw=black,very thick,minimum size=1.9cm,font=\small,align=center}]
  \node[ev,fill=gray!8] (u) {up-move\\(type 1)};
  \node[ev,fill=gray!25,right=of u] (d) {down-move\\(type 2)};
  \draw[-{Stealth[length=2.5mm]},very thick] (u) edge[loop left,looseness=6]
    node[left]{$\alpha_{\mathrm{self}}$} (u);
  \draw[-{Stealth[length=2.5mm]},very thick] (d) edge[loop right,looseness=6]
    node[right]{$\alpha_{\mathrm{self}}$} (d);
  \draw[-{Stealth[length=2.5mm]},thick,gray!70!black] (u) to[bend left=22]
    node[above]{$\alpha_{\mathrm{cross}}$} (d);
  \draw[-{Stealth[length=2.5mm]},thick,gray!70!black] (d) to[bend left=22]
    node[below]{$\alpha_{\mathrm{cross}}$} (u);
  \node[font=\small] at ($(u)!0.5!(d)+(0,-1.9)$) {baseline $\mu$ for each type; all excitation decays at rate $\beta$};
\end{tikzpicture}
\caption{The symmetric bivariate Hawkes process as a diagram of event types. Loops are
self-excitation, arrows between types are cross-excitation; the same four parameters
\((\mu,\alpha_{\mathrm{self}},\alpha_{\mathrm{cross}},\beta)\) apply to both types.
\(\alpha_{\mathrm{self}}>\alpha_{\mathrm{cross}}\) describes persistence of direction, as
found in the crypto data of this work; \(\alpha_{\mathrm{cross}}>\alpha_{\mathrm{self}}\)
describes reversion, as at the equity fixed point of \citet{lee2023}.}
\label{fig:bivariate}
\end{figure}

\section{Stability}

The mass matrix of the kernels is
\[
 \mathbf G=\frac{1}{\beta}
 \begin{bmatrix}
  \alpha_{\mathrm{self}} & \alpha_{\mathrm{cross}}\\
  \alpha_{\mathrm{cross}} & \alpha_{\mathrm{self}}
 \end{bmatrix}.
\]
Its entry \((i,j)\) is the mean number of type-\(i\) children of a type-\(j\) event, so
\(\mathbf G\) plays the role of the branching ratio. Its eigenvalues are
\((\alpha_{\mathrm{self}}-\alpha_{\mathrm{cross}})/\beta\) and the largest,
\[
 g=\frac{\alpha_{\mathrm{self}}+\alpha_{\mathrm{cross}}}{\beta},
\]
which is the mean number of children of either type per event. The two eigenvalues have
eigenvectors \((1,1)\) and \((1,-1)\): the first governs the total activity
\(N_1+N_2\), the second the difference \(N_1-N_2\), that is, the price
(Section~\ref{sec:volatility-theory}). Write
\[
 g_-=\frac{\alpha_{\mathrm{self}}-\alpha_{\mathrm{cross}}}{\beta}
\]
for the second eigenvalue. The condition \(g<1\) is
required. Close to one, cascades last longer; above one, there is no usual stationary
regime \citep{hardiman2013,filimonov2012}. Under stationarity, by the same argument as in
the univariate case, each type occurs at mean rate \(\mu/(1-g)\), and the two together at
\(2\mu/(1-g)\).

\section{Marks and price}

An event has a time, a type and a mark. The mark is the size of the change in ticks.
The paper keeps the mark for volatility, but assumes that the intensity does not
depend on it. The filtered price can be written, in ticks, as the cumulative
difference between up-moves and down-moves.

\section{Volatility}
\label{sec:volatility-theory}

In a minimal model, if the price moves by one tick \(\vartheta\) at a rate \(\phi\),
\[
 \sigma^2\approx\phi\vartheta^2.
\]
Thus, modeling intensity is modeling activity and, consequently, a component of
volatility. The marked Hawkes expression requires moments of the marks and the
estimated parameters \citep{leeseo2017,lee2024}.

\paragraph{Long-run variance of the Hawkes price.} With unit marks the price is
\(X(t)=N_1(t)-N_2(t)\), the projection on the eigenvector \((1,-1)\). Under stationarity each
type has mean rate \(\bar\lambda=\mu/(1-g)\), and the variance of \(X\) per unit time
converges, as the horizon grows, to
\[
 \lim_{T\to\infty}\frac{\Var X(T)}{T}=\frac{2\bar\lambda}{(1-g_-)^2}
\]
\citep{bacry2013}. If each event moves the price by an i.i.d.\ size \(J\), independent of the
times, with the sign given by the type, the same argument gives
\begin{equation}
 \sigma_H^2=2\bar\lambda\left[\Var(J)+\frac{\E[J]^2}{(1-g_-)^2}\right],
 \label{eq:hawkes-variance}
\end{equation}
the expression used in Section~\ref{sec:real-volatility}. The total activity depends on
\(g\), and the amplification of the variance on \(g_-\). When self-excitation dominates,
\(g_->0\) and moves in the same direction cluster, so the variance exceeds that of a Poisson
price with the same activity. When cross-excitation dominates, \(g_-<0\), moves tend to undo
each other and the variance falls below it; this is the mechanism that \citet{bacry2013} use
to model microstructure noise. When \(\alpha_{\mathrm{self}}=\alpha_{\mathrm{cross}}\), the
long-run variance equals that of the difference of two independent Poisson processes with the
same activity.

\paragraph{Realized volatility.} The benchmark is the realized variance over a window,
\[
 \mathrm{RV}_\Delta=\sum_k r_k^2,
\]
the sum of squared log returns sampled every \(\Delta\) seconds
\citep{andersen2003,barndorff2002}. Without noise, it converges to the integrated variance as
\(\Delta\to0\). At high frequency, however, the observed price contains microstructure noise,
and \(\mathrm{RV}_\Delta\) is biased upward when \(\Delta\) is small \citep{zhang2005}. The
signature plot, \(\mathrm{RV}_\Delta\) per unit time against \(\Delta\), shows the effect:
noise makes it fall with \(\Delta\), while positively autocorrelated returns, as in a
self-exciting price, make it rise toward the long-run level of
Equation~\eqref{eq:hawkes-variance}.

\section{Simulation}

A Hawkes path can be simulated by Ogata's thinning \citep{ogata1981}. If
\(\lambda(t)\le M\) on an interval, candidates are drawn from a Poisson process of rate
\(M\) and each candidate at time \(s\) is kept with probability \(\lambda(s)/M\); the kept
points have intensity \(\lambda\). With an exponential kernel the intensity only decreases
between events, so its value just after the last event is a valid bound until the next
event. In the bivariate case the bound is \(\lambda_1+\lambda_2\), and the type of an
accepted event is drawn in proportion to \(\lambda_1\) and \(\lambda_2\). The
implementation is described in Section~\ref{sec:sim-thinning}.

\section{Maximum likelihood}
\label{sec:mle-theory}

For a horizon \(T\),
\[
 \ell(\theta)=\sum_k\log\lambda_{c_k}(t_k)
 -\sum_i\int_0^T\lambda_i(u)\,du.
\]
The first term rewards high intensity at the events; the second penalizes high
intensity over the whole interval. Maximizing only the first would drive the
intensities to infinity. The compensator provides the balance. For an exponential
kernel, the sum over the past obeys a first-order recursion, and \(\ell\) and its
gradient can be computed in \(O(n)\) operations \citep{ozaki1979}. In the univariate case,
let \(A_1=0\) and \(A_k=e^{-\beta(t_k-t_{k-1})}(1+A_{k-1})\), so that \(\alpha A_k\) is
the excitation just before \(t_k\) (Equation~\eqref{eq:recursion}). Then
\begin{equation}
 \ell(\mu,\alpha,\beta)=\sum_{k=1}^{n}\log\bigl(\mu+\alpha A_k\bigr)-\mu T
 -\frac{\alpha}{\beta}\sum_{k=1}^{n}\Bigl(1-e^{-\beta(T-t_k)}\Bigr).
 \label{eq:loglik1d}
\end{equation}
The bivariate likelihood has the same structure, with one recursion per pair of source and
target types. Under regularity conditions, the MLE is consistent and asymptotically
efficient \citep{ogata1978}: approximately normal around the true value, with covariance
given by the inverse Fisher information, which the inverse of the observed Hessian of
\(-\ell\) estimates. These conditions include a true parameter in the interior of the
parameter space. When it lies on the boundary, as \(\alpha_{\mathrm{cross}}=0\), the
estimate piles up at the boundary with positive probability, and neither the normal
approximation nor the usual \(\chi^2\) reference for likelihood ratio tests applies
\citep{selfliang1987}. This is relevant for the real data of this work, where
\(\alpha_{\mathrm{cross}}\) is estimated close to zero.

\section{Model comparison}

The Poisson process is nested in the Hawkes process: it is the case
\(\alpha_{\mathrm{self}}=\alpha_{\mathrm{cross}}=0\). The likelihood ratio statistic
\(2(\hat\ell_{\mathrm{Hawkes}}-\hat\ell_{\mathrm{Poisson}})\) measures the gain in fit. Its
usual \(\chi^2\) reference does not hold here for two reasons: the null value lies on the
boundary of the parameter space \citep{selfliang1987}, and \(\beta\) has no effect under the
null, so it is not identified there. The Akaike information criterion,
\[
 \mathrm{AIC}=2k-2\hat\ell,
\]
with \(k\) the number of free parameters, compares the two fits with a penalty for the extra
parameters \citep{akaike1974}. A difference in AIC of hundreds of units, as found on the real
windows, leaves no doubt about the direction, whatever the exact reference distribution.

\section{Goodness of fit by time rescaling}
\label{sec:time-rescaling}

If the model is correct, the compensator turns the observed process into a unit-rate
Poisson process: the transformed times \(\Lambda(t_k)\) are the events of a homogeneous
Poisson process of rate one. This is the time-rescaling theorem \citep{brown2002}. The
increments
\[
 \tau_k=\Lambda(t_k)-\Lambda(t_{k-1})=\int_{t_{k-1}}^{t_k}\lambda(u)\,\mathrm du
\]
are then independent \(\mathrm{Exp}(1)\) variables. In the bivariate case the theorem
applies to each type with its own compensator \(\Lambda_i\). Evaluated at the estimated
parameters, the \(\tau_k\) are the rescaled residuals of the fit \citep{ogata1988}. Their
distribution is tested against \(\mathrm{Exp}(1)\), for example with Kolmogorov--Smirnov,
and their serial dependence with Ljung--Box. A kernel of the wrong shape typically leaves
dependence in the residuals: bursts that the fitted kernel forgets too fast or too slowly
appear as runs of short or long \(\tau_k\).

\section{Observation on a time grid}

The events of this work are changes of the mid-price recorded on a 100~ms grid
(Chapter~\ref{cap:dados}). Two things are lost relative to continuous time: event times are
rounded to the grid, and changes that fall in the same cell merge into one. The second
matters most for the decay. With \(\beta\) of a few units per second, much of the excitation
acts within a few hundred milliseconds, and merged events remove part of the short
inter-arrivals from which \(\beta\) is identified. A likelihood written for continuous time
therefore estimates the parameters of the process as seen through the grid, with \(\beta\)
biased downward; Section~\ref{sec:addendum-grid} measures the size of this bias.
Likelihoods that account for the observation scheme have been proposed for this setting
\citep{lambe2025}.

A biased estimator can still be corrected by simulation, the idea of indirect inference
\citep{gourieroux1993}. Let \(b(\theta)\) be a summary, such as the median, of the grid MLE
over paths simulated from \(\theta\) in continuous time and mapped to the grid. The map
\(b\) is estimated by simulation, and an estimate on the grid is corrected by solving
\(b(\theta)=\hat\theta_{\mathrm{grid}}\). The correction is as good as the model used to
simulate: if the real process is not a Hawkes process, the map transfers its bias with it.
A network trained on gridded paths with the continuous-time parameter as target learns a
correction of the same kind, conditional on the whole path.

\section{Identifiability in finite samples}
\label{sec:identificabilidade}

The asymptotic properties do not guarantee good behavior with 1,000 events. When the
excitation is weak (small \(\alpha_{\mathrm{self}}+\alpha_{\mathrm{cross}}\)), the
kernel barely changes the intensity, and the likelihood becomes flat in the direction
of \(\beta\): very different values of \(\beta\) explain the data almost equally well.
In the limit \(\alpha\to0\), \(\beta\) is no longer identified. Near criticality, the
separation between \(\mu\) and \(g\) also deteriorates, because the mean rate
\(2\mu/(1-g)\) is well determined, but neither factor is on its own. These effects
reappear in the results of Chapter~\ref{cap:resultados}.

\section{Amortized estimation as a Bayes estimator}
\label{sec:bayes}

In the approach of \citet{lee2023}, a network \(f_\phi\) is trained on pairs
\((\theta_i,x_i)\), where \(\theta_i\) is drawn from a training distribution \(\pi\)
and \(x_i\) is simulated from the model with parameter \(\theta_i\). Minimizing the
mean squared error,
\[
  \min_\phi\;\E_{\theta\sim\pi}\,\E_{x\mid\theta}
  \bigl\lVert f_\phi(x)-\theta\bigr\rVert^2,
\]
is minimizing the Bayes risk under squared loss, whose solution is the posterior mean
\(\E_\pi[\theta\mid x]\) \citep{lehmann1998}. Hence, a well-trained network
approximates a Bayes estimator under the prior \(\pi\)
\citep{sainsburydale2024,lenzi2023}. Two consequences guide this work. First, on a
test set drawn from the same \(\pi\), no estimator has lower average MSE than the
posterior mean, so outperforming the MLE on that test is expected and does not prove
superiority in a frequentist sense. Second, the network does not extrapolate: its
estimates shrink toward the high-density region of \(\pi\), and real data outside the
support of the prior yield biased estimates.

The same approach belongs to a wider family of amortized, simulation-based inference methods,
in which the cost of inference is paid once, in training, and each new data set costs a
forward pass \citep{cranmer2020,radev2022}. On real windows the target changes. When a
network is trained on pairs \((x,\hat\theta_{\mathrm{MLE}}(x))\), the label is a deterministic
function of the input, so the minimizer of the squared loss is the MLE itself: the network
becomes a fast approximation of the MLE map, accurate where the training windows are dense.
Starting from the weights trained on simulated data, it can reach that approximation with
fewer real windows; how close the simulation prior is to the real data then matters for how
much that initialization helps.

\section{Comparing estimators}
\label{sec:comparing-theory}

The estimators are compared path by path. For each path or window \(i\), a loss
\(L_i(A)\), such as the squared error averaged over the four parameters, is computed for each
estimator \(A\); the macro MSE is the average of the per-parameter MSEs. Two estimators are
compared through the paired differences \(d_i=L_i(A)-L_i(B)\). Pairing removes the variation
in difficulty between paths, which is often much larger than the difference between
estimators.
The mean \(\bar d\) is reported with a bootstrap confidence interval, obtained by resampling
the \(d_i\) with replacement \citep{efron1993}, and the null hypothesis of no difference is
tested with the Wilcoxon signed-rank test, which uses the ranks of \(|d_i|\) and is robust to
the heavy tails of squared errors \citep{wilcoxon1945,diebold1995}. Variation across training
seeds is a different source: it measures the instability of training, not the accuracy of a
trained model, and is compared across seeds \citep{bouthillier2021}.

When several comparisons are tested together, the Holm procedure controls the probability
of any false rejection: with the \(m\) \(p\)-values sorted in increasing order, the \(k\)-th
is compared with \(0.05/(m-k+1)\), and testing stops at the first non-rejection
\citep{holm1979}. Agreement between two orderings, such as the Hawkes and the realized
volatilities across windows, is measured by the Spearman correlation, the correlation of the
ranks, which ignores the level and depends only on the order.

\paragraph{Ensembles.} The average \(\bar y=\frac1M\sum_m y_m\) of \(M\) estimators satisfies
\[
 \lVert\bar y-\theta\rVert^2=\frac1M\sum_{m=1}^{M}\lVert y_m-\theta\rVert^2
 -\frac1M\sum_{m=1}^{M}\lVert y_m-\bar y\rVert^2
\]
\citep{krogh1995}. The error of the average equals the average error of the members minus
their spread, so it is never worse than the average member, and it improves most when the
members are accurate and err in different places. This identity applies both to the
ensembles of five seeds and to the average of the LSTM and the Transformer.

\section{LSTM}

Recurrent networks update a state \(h_t=f(x_t,h_{t-1})\). The LSTM introduces input,
forget and output gates to control the flow of information and mitigate the
vanishing gradient problem \citep{hochreiter1997}. With \(\sigma\) the logistic function
and \(\odot\) the elementwise product,
\begin{align*}
 f_t&=\sigma(W_f[h_{t-1},x_t]+b_f), & i_t&=\sigma(W_i[h_{t-1},x_t]+b_i),\\
 o_t&=\sigma(W_o[h_{t-1},x_t]+b_o), & \tilde c_t&=\tanh(W_c[h_{t-1},x_t]+b_c),\\
 c_t&=f_t\odot c_{t-1}+i_t\odot\tilde c_t, & h_t&=o_t\odot\tanh(c_t).
\end{align*}
The cell state \(c_t\) is updated additively, which lets gradients pass over many steps.
The update has the form of Equation~\eqref{eq:recursion}: a forget gate that depends on the
inter-arrival time can play the role of \(e^{-\beta\Delta_k}\), and the input gate that of
the jump, so an LSTM can represent the excitation of a Hawkes process in its state. In
Lee, the sequence contains the inter-arrival time and the type; the final state feeds the
regression of the four parameters. The paper also tested a GRU \citep{cho2014}, with
slightly worse results.

\section{Attention and the Transformer}

Attention computes weighted combinations:
\[
 \operatorname{Attention}(Q,K,V)=
 \operatorname{softmax}\left(\frac{QK^\top}{\sqrt{d_k}}\right)V.
\]
Each position queries all the others. The Transformer combines multiple heads,
feed-forward layers, residual connections and normalization \citep{vaswani2017}. Each head
applies its own projections \(Q=XW^Q\), \(K=XW^K\), \(V=XW^V\) to the sequence \(X\), and
the heads are concatenated. Self-attention is equivariant to permutations of the sequence,
so order and spacing must enter through a positional encoding; for a point process the
relevant position is the time of each event, not only its index. The
potential advantage is capturing distant dependencies; the potential cost is
\(O(n^2)\) memory. The position of the normalization relative to the residual
connections affects training stability: the \emph{post-LN} variant, used here,
usually requires learning-rate warm-up \citep{xiong2020}.

\section{Training}

The networks are trained by stochastic gradient descent on mini-batches with Adam, which
scales each step by running estimates of the first and second moments of the gradient
\citep{kingma2015}. The learning rate can follow a schedule, such as a short warm-up or a cosine decay
\citep{xiong2020,loshchilov2017}. Training stops at a fixed budget of epochs or when the
validation loss stops improving, and the checkpoint with the lowest validation loss is kept;
the test data are not used for any choice. Fine-tuning reuses this procedure from weights
already trained on another task, here the simulated paths, with the real windows as the new
training set.
